% Generated by source/build_book.py from the \keyterm definitions. Do not edit by hand.
\section*{Glossary}
\addcontentsline{toc}{section}{Glossary}
\begingroup\small\setlength{\parindent}{-1.5em}\setlength{\leftskip}{1.5em}\RaggedRight
\textbf{appreciation}\quad A rise in the value of a currency set by the market; here, fewer rupees are needed to buy one dollar.\par\smallskip
\textbf{authorised dealers}\quad Banks and firms licensed by the RBI under FEMA, 1999, to buy and sell foreign currency.\par\smallskip
\textbf{balance of payments (BoP)}\quad The record of all economic transactions between a country's residents and the rest of the world over a period.\par\smallskip
\textbf{capital account}\quad The part of the BoP that records the buying and selling of assets, such as FDI, FPI and loans.\par\smallskip
\textbf{capital flows}\quad Money that moves between countries to buy or sell financial assets, such as shares and bonds, or real assets, such as factories.\par\smallskip
\textbf{currency crisis}\quad A sudden and large fall in the value of a country's currency, usually because the country has run short of foreign exchange.\par\smallskip
\textbf{current account}\quad The part of the BoP that records trade in goods and services, income and transfers.\par\smallskip
\textbf{current account deficit (CAD)}\quad The gap when a country's payments abroad for goods, services, income and transfers exceed its receipts.\par\smallskip
\textbf{demand for foreign exchange}\quad The foreign currency that residents want, to pay for imports, study and travel abroad, investment abroad and foreign loans.\par\smallskip
\textbf{depreciation}\quad A fall in the value of a currency set by the market; here, more rupees are needed to buy one dollar.\par\smallskip
\textbf{devaluation}\quad A deliberate cut in a currency's official value by the government or central bank under a fixed exchange rate.\par\smallskip
\textbf{direct quote}\quad The price of one unit of foreign currency in home currency, e.g.\ US\$1 = ₹95.\par\smallskip
\textbf{equilibrium exchange rate}\quad The exchange rate at which the quantity of dollars demanded equals the quantity supplied.\par\smallskip
\textbf{excess demand}\quad When buyers want more dollars than are on offer at the current exchange rate; the rate tends to rise.\par\smallskip
\textbf{excess supply}\quad When more dollars are offered than buyers want at the current exchange rate; the rate tends to fall.\par\smallskip
\textbf{exchange rate}\quad The price of one currency in terms of another, for example ₹95 per US dollar.\par\smallskip
\textbf{exchange rate channel}\quad The way a change in the central bank's policy interest rate affects spending and prices by moving the exchange rate, and through it exports and imports.\par\smallskip
\textbf{exchange rate system}\quad The rules a country follows to manage the value of its currency against other currencies.\par\smallskip
\textbf{export competitiveness}\quad The ability of a country's firms to sell goods and services abroad at prices that foreign buyers are willing to pay.\par\smallskip
\textbf{fixed (pegged) exchange rate}\quad An exchange rate held at a set value against another currency, a basket of currencies or gold by the central bank.\par\smallskip
\textbf{floating (flexible) exchange rate}\quad An exchange rate set by demand and supply in the market, with little government control.\par\smallskip
\textbf{foreign direct investment (FDI)}\quad Buying a lasting stake (10\% or more) in a business abroad, or setting one up.\par\smallskip
\textbf{foreign exchange market}\quad The market in which currencies are traded for one another. Also called the forex market.\par\smallskip
\textbf{foreign exchange reserves}\quad Foreign currency assets, gold and other reserve assets held by the RBI on behalf of the country.\par\smallskip
\textbf{foreign portfolio investment (FPI)}\quad Buying shares and bonds (holdings under 10\%), which can be sold quickly.\par\smallskip
\textbf{iMF programme}\quad A loan from the International Monetary Fund to a country short of foreign currency, paid in instalments on condition that the country carries out agreed policy changes.\par\smallskip
\textbf{import cover}\quad The number of months of imports a country could pay for using only its reserves.\par\smallskip
\textbf{imported inflation}\quad A rise in the general price level caused by higher rupee prices of imported goods, often because the rupee has depreciated.\par\smallskip
\textbf{indirect quote}\quad How much foreign currency one unit of home currency buys, e.g.\ ₹1 = US\$0.0105.\par\smallskip
\textbf{interest rate differential}\quad The difference between the interest rates paid in two countries. Funds tend to move towards the higher return, after allowing for risk.\par\smallskip
\textbf{intervention}\quad Buying or selling foreign currency by the central bank to influence the exchange rate.\par\smallskip
\textbf{invisibles}\quad Trade in services, income and transfers, as opposed to trade in goods.\par\smallskip
\textbf{managed float}\quad An exchange rate mainly set by the market, with the central bank stepping in to prevent sharp swings.\par\smallskip
\textbf{nominal exchange rate}\quad The number of rupees needed to buy one unit of foreign currency, as quoted in the market.\par\smallskip
\textbf{purchasing power parity (PPP)}\quad The idea that the same basket of goods should cost the same in every country once prices are converted into one currency.\par\smallskip
\textbf{real exchange rate}\quad The nominal exchange rate adjusted for differences in price levels between two countries.\par\smallskip
\textbf{remittances}\quad Money that people working abroad send home to their families.\par\smallskip
\textbf{revaluation}\quad A deliberate increase in a currency's official value under a fixed exchange rate.\par\smallskip
\textbf{sovereign default}\quad When a government fails to make the payments on its debt that it has promised.\par\smallskip
\textbf{Special Drawing Rights (SDR)}\quad An international reserve asset created by the IMF that member countries can use to settle payments with each other.\par\smallskip
\textbf{speculation}\quad Buying or selling an asset, such as a currency, in the hope of profiting from a change in its price.\par\smallskip
\textbf{supply of foreign exchange}\quad The foreign currency flowing into a country from its exports, remittances, tourists and foreign investment.\par\smallskip
\endgroup
